\documentclass[12pt]{article}

\usepackage{amsmath}
\usepackage{amssymb}
\usepackage{physics}
\usepackage{hyperref}

\title{Hyperplanes and Linear Equations}
\author{}
\date{}

\begin{document}

\maketitle

\section{Analytic Geometry}

A functional on a real vector space \(V\) is a mapping
\[
f:V\longrightarrow\mathbb{R}
\]
which satisfies the linearity conditions
\[
f(\mathbf{u}+\mathbf{v})
=
f(\mathbf{u})+f(\mathbf{v}),
\qquad
\mathbf{u},\mathbf{v}\in V,
\]
and
\[
f(t\mathbf{v})
=
t f(\mathbf{v}),
\qquad
\mathbf{v}\in V,\quad t\in\mathbb{R}.
\]

\begin{proposition}
Let \(V\) be a real vector space of dimension \(n\). Any
\((n-1)\)-dimensional linear subspace \(U\) of \(V\) is the kernel of a
nonzero functional
\[
f:V\longrightarrow\mathbb{R}.
\]
Conversely, the kernel of every nonzero functional on \(V\) is an
\((n-1)\)-dimensional linear subspace of \(V\). Moreover, the functional
is determined by \(U\) up to a nonzero scalar factor.
\end{proposition}

\begin{proof}
We first prove that the kernel of a nonzero functional \(f\) is a linear
subspace of \(V\) of dimension \(n-1\).

The kernel of \(f\) is
\[
\ker(f)
=
\left\{
\mathbf{u}\in V
\mid
f(\mathbf{u})=0
\right\}.
\]
Let \(\mathbf{u},\mathbf{v}\in\ker(f)\), and let
\(\alpha,\beta\in\mathbb{R}\). By linearity,
\[
f(\alpha\mathbf{u}+\beta\mathbf{v})
=
\alpha f(\mathbf{u})+\beta f(\mathbf{v})
=
0.
\]
Hence
\[
\alpha\mathbf{u}+\beta\mathbf{v}\in\ker(f),
\]
and therefore \(\ker(f)\) is a linear subspace of \(V\).

Let
\[
\mathbf{v}_1,\ldots,\mathbf{v}_n
\]
be a basis of \(V\). Since \(f\) is nonzero, at least one basis vector
has a nonzero image. After relabeling the basis if necessary, we may
assume that
\[
f(\mathbf{v}_n)\neq0.
\]
For any \(\mathbf{u}\in\ker(f)\), write
\[
\mathbf{u}=u^i\mathbf{v}_i.
\]
Then
\[
f(\mathbf{u})
=
u^i f(\mathbf{v}_i)
=
0.
\]
Thus
\[
u^n
=
-\frac{
u^1f(\mathbf{v}_1)+\cdots+u^{n-1}f(\mathbf{v}_{n-1})
}{
f(\mathbf{v}_n)
}.
\]
Therefore only \(n-1\) coefficients can be chosen independently, and
\[
\dim\ker(f)=n-1.
\]

We now prove the converse. Let \(U\) be an \((n-1)\)-dimensional linear
subspace of \(V\). Choose a vector
\[
\mathbf{e}\in V\setminus U.
\]
Then
\[
W=\operatorname{span}\{\mathbf{e}\}
\]
is a one-dimensional linear subspace of \(V\), and
\[
U\cap W=\{\mathbf{0}\}.
\]
Since
\[
\dim U+\dim W=n,
\]
we have
\[
V=U\oplus W.
\]
Hence every \(\mathbf{v}\in V\) can be written uniquely as
\[
\mathbf{v}=\mathbf{u}+\beta\mathbf{e},
\qquad
\mathbf{u}\in U,\quad \beta\in\mathbb{R}.
\]

Define
\[
f:V\longrightarrow\mathbb{R}
\]
by
\[
f(\mathbf{v})=\beta.
\]
This functional is linear and nonzero. Moreover,
\[
f(\mathbf{v})=0
\]
if and only if \(\beta=0\), which holds if and only if
\(\mathbf{v}\in U\). Therefore,
\[
\ker(f)=U.
\]

Finally, let \(f\) and \(g\) be two nonzero functionals with the same
kernel \(U\):
\[
\ker(f)=\ker(g)=U.
\]
For any \(\mathbf{v}\in V\), write
\[
\mathbf{v}=\mathbf{u}+\beta\mathbf{e},
\qquad
\mathbf{u}\in U.
\]
Since
\[
f(\mathbf{u})=g(\mathbf{u})=0,
\]
we have
\[
f(\mathbf{v})=\beta f(\mathbf{e})
\]
and
\[
g(\mathbf{v})=\beta g(\mathbf{e}).
\]
Furthermore,
\[
f(\mathbf{e})\neq0
\qquad\text{and}\qquad
g(\mathbf{e})\neq0,
\]
because \(\mathbf{e}\notin U\). Define
\[
\lambda
=
\frac{f(\mathbf{e})}{g(\mathbf{e})}.
\]
Then \(\lambda\neq0\), and
\[
f(\mathbf{v})
=
\lambda g(\mathbf{v})
\]
for every \(\mathbf{v}\in V\). Hence \(f\) is determined by \(U\) up to
a nonzero scalar factor.
\end{proof}

\begin{definition}
Let \(\mathcal A\) be an \(n\)-dimensional affine space. An affine
subspace \(S(x,U)\) of \(\mathcal A\) is called a \emph{hyperplane} if its
direction space \(U\) has dimension \(n-1\).
\end{definition}

\begin{definition}[Coordinate map]
Let
\[
x_0,\mathbf v_1,\ldots,\mathbf v_n
\]
be an affine coordinate system for an affine space $(\mathcal A,V)$, where
$\mathbf v_1,\ldots,\mathbf v_n$ is a basis of $V$.

The coordinate map associated with this affine coordinate system is the map
\[
c:\mathcal A\longrightarrow\mathbb R^n
\]
defined by
\[
c(y)=(y^1,\ldots,y^n)
\]
if and only if
\[
y=(y^i\mathbf v_i)(x_0).
\]
\end{definition}

\begin{proposition}
Let $(\mathcal A,V)$ be an affine space of dimension $n$, equipped with the
affine coordinate system
\[
x_0,\mathbf v_1,\ldots,\mathbf v_n.
\]

A hyperplane of $\mathcal A$ is determined, with respect to this reference
system, by a linear equation
\[
d^i\xi_i=\mu_0,
\]
where $\xi_1,\ldots,\xi_n,\mu_0\in\mathbb R$ and not all the
coefficients $\xi_i$ are zero.

Conversely, every equation of this form determines a hyperplane of
$\mathcal A$.
\end{proposition}

\begin{proof}
Let us first consider a hyperplane
\[
S(x,U)=\{\mathbf u(x)\mid \mathbf u\in U\},
\]
where $U$ is an $(n-1)$-dimensional linear subspace of $V$.

By the preceding proposition, there exists a nonzero functional
\[
f:V\longrightarrow\mathbb R
\]
such that
\[
\ker f=U.
\]

Define the set of coordinate tuples of the points of the hyperplane by
\[
C:=c\bigl(S(x,U)\bigr)
  =\left\{(d^i)\in\mathbb R^n
    \,\middle|\,
    \exists\,y\in S(x,U),\ c(y)=(d^i)
   \right\}.
\]

Because the action of $V$ on $\mathcal A$ is transitive, there exists a unique
vector $\mathbf u_0\in V$ such that
\[
x=\mathbf u_0(x_0).
\]

We shall prove that
\[
C=F,
\]
where
\[
F:=\left\{(d^i)\in\mathbb R^n
          \,\middle|\,
          d^i\xi_i=\mu_0
   \right\}
\]
for suitable coefficients $\xi_i$ and $\mu_0$.

Let $(d^i)\in C$. Then there exists a point
\[
y\in S(x,U)
\]
whose coordinates are $(d^i)$. Define its displacement vector by
\[
\mathbf d:=\overrightarrow{x_0y}=d^i\mathbf v_i,
\qquad
y=\mathbf d(x_0).
\]

Since $y\in S(x,U)$, there exists a vector $\mathbf u\in U$ such that
\[
y=\mathbf u(x).
\]
Using $x=\mathbf u_0(x_0)$, we obtain
\[
y
 =\mathbf u\bigl(\mathbf u_0(x_0)\bigr)
 =(\mathbf u+\mathbf u_0)(x_0).
\]
On the other hand,
\[
y=\mathbf d(x_0).
\]
Since the action is free,
\[
\mathbf d=\mathbf u+\mathbf u_0,
\]
and therefore
\[
\mathbf u=\mathbf d-\mathbf u_0.
\]

Because $\mathbf u\in U=\ker f$,
\[
f(\mathbf d-\mathbf u_0)=0.
\]
By the linearity of $f$,
\[
f(\mathbf d)-f(\mathbf u_0)=0.
\]
Since $\mathbf d=d^i\mathbf v_i$, this becomes
\[
d^if(\mathbf v_i)-f(\mathbf u_0)=0.
\]

We now define
\[
\xi_i:=f(\mathbf v_i),
\qquad
\mu_0:=f(\mathbf u_0).
\]
It follows that
\[
d^i\xi_i=\mu_0.
\]
Therefore
\[
(d^i)\in F,
\]
and hence
\[
C\subseteq F.
\]

Moreover, not all the coefficients $\xi_i$ are zero. Indeed, if
\[
\xi_i=f(\mathbf v_i)=0
\qquad
\text{for every }i,
\]
then $f$ would vanish on a basis of $V$ and would therefore be the zero
functional, contrary to the choice of $f$.

We now prove the reverse inclusion. Let
\[
(d^i)\in F
\]
and define
\[
\mathbf d:=d^i\mathbf v_i,
\qquad
y:=\mathbf d(x_0).
\]
Because $(d^i)\in F$,
\[
d^i\xi_i=\mu_0.
\]
Consequently,
\[
f(\mathbf d-\mathbf u_0)
 =d^if(\mathbf v_i)-f(\mathbf u_0)
 =d^i\xi_i-\mu_0
 =0.
\]
Thus
\[
\mathbf d-\mathbf u_0\in\ker f=U.
\]

Since $x=\mathbf u_0(x_0)$, we have
\[
y
 =\mathbf d(x_0)
 =(\mathbf d-\mathbf u_0)\bigl(\mathbf u_0(x_0)\bigr)
 =(\mathbf d-\mathbf u_0)(x).
\]
Therefore
\[
y\in S(x,U),
\]
and hence
\[
(d^i)=c(y)\in C.
\]
It follows that
\[
F\subseteq C.
\]

We conclude that
\[
C=F.
\]
Thus every hyperplane is determined, in the chosen affine coordinate system,
by one linear equation
\[
d^i\xi_i=\mu_0,
\]
where not all the coefficients $\xi_i$ are zero.

We now prove the converse. Consider the set
\[
F:=\left\{(d^i)\in\mathbb R^n
          \,\middle|\,
          d^i\xi_i=\mu_0
   \right\},
\]
where not all the coefficients $\xi_i$ are zero.

Define a functional
\[
f:V\longrightarrow\mathbb R
\]
by prescribing its values on the basis:
\[
f(\mathbf v_i):=\xi_i.
\]
Equivalently, if
\[
\mathbf d=d^i\mathbf v_i,
\]
then
\[
f(\mathbf d)=d^i\xi_i.
\]

Since not all the coefficients $\xi_i$ are zero, the functional $f$ is
nonzero. Therefore
\[
U:=\ker f
\]
is an $(n-1)$-dimensional linear subspace of $V$.

Choose a vector
\[
\mathbf m\in V\setminus U
\]
and define the one-dimensional subspace
\[
W:=\operatorname{span}\{\mathbf m\}.
\]
Since $\mathbf m\notin\ker f$, the restriction
\[
f|_W:W\longrightarrow\mathbb R
\]
is nonzero. Because both $W$ and $\mathbb R$ are one-dimensional,
$f|_W$ is surjective. Hence there exists a vector
\[
\mathbf v_0\in W
\]
such that
\[
f(\mathbf v_0)=\mu_0.
\]

Define the point
\[
x:=\mathbf v_0(x_0)
\]
and consider the affine subspace
\[
S(x,U).
\]
Since $\dim U=n-1$, this affine subspace is a hyperplane.

Let
\[
C:=c\bigl(S(x,U)\bigr).
\]
We shall prove that
\[
C=F.
\]

Let $(d^i)\in C$. Then there exists a point
\[
y\in S(x,U)
\]
with
\[
c(y)=(d^i).
\]
Define
\[
\mathbf d:=d^i\mathbf v_i,
\]
so that
\[
y=\mathbf d(x_0).
\]

Since $y\in S(x,U)$, there exists $\mathbf u\in U$ such that
\[
y=\mathbf u(x).
\]
Using $x=\mathbf v_0(x_0)$, we obtain
\[
y
 =\mathbf u\bigl(\mathbf v_0(x_0)\bigr)
 =(\mathbf u+\mathbf v_0)(x_0).
\]
Since also $y=\mathbf d(x_0)$ and the action is free,
\[
\mathbf d=\mathbf u+\mathbf v_0.
\]
Therefore
\[
f(\mathbf d)
 =f(\mathbf u)+f(\mathbf v_0)
 =0+\mu_0
 =\mu_0.
\]
Since
\[
f(\mathbf d)=d^i\xi_i,
\]
we obtain
\[
d^i\xi_i=\mu_0.
\]
Thus
\[
(d^i)\in F,
\]
and hence
\[
C\subseteq F.
\]

Conversely, let
\[
(d^i)\in F
\]
and define
\[
\mathbf d:=d^i\mathbf v_i,
\qquad
y:=\mathbf d(x_0).
\]
Since $(d^i)\in F$,
\[
d^i\xi_i=\mu_0.
\]
It follows that
\[
f(\mathbf d-\mathbf v_0)
 =f(\mathbf d)-f(\mathbf v_0)
 =d^i\xi_i-\mu_0
 =0.
\]
Therefore
\[
\mathbf d-\mathbf v_0\in\ker f=U.
\]

Using $x=\mathbf v_0(x_0)$, we obtain
\[
y
 =\mathbf d(x_0)
 =(\mathbf d-\mathbf v_0)\bigl(\mathbf v_0(x_0)\bigr)
 =(\mathbf d-\mathbf v_0)(x).
\]
Hence
\[
y\in S(x,U),
\]
so
\[
(d^i)=c(y)\in C.
\]
Therefore
\[
F\subseteq C.
\]

We conclude that
\[
C=F.
\]
Thus the given linear equation determines the hyperplane $S(x,U)$.
\end{proof}

\end{document}
